\subsection{Construction of polynomials}

Let A be the ring Z{[}X, Y{]} of polynomials with integer coefficients in the two families of indeterminates \(\mathbf{X} = (\mathbf{X}_{n})_{n \in \mathbb{N}}\) and \(\mathbf{Y} = (\mathbf{Y}_{n})_{n \in \mathbb{N}}\). Let \(\theta\) be the endomorphism of A defined by \(\theta(\mathbf{X}_{n}) = \mathbf{X}_{n}^{p}\) and \(\theta(\mathbf{Y}_{n}) = \mathbf{Y}_{n}^{p}\) for every \(n \in \mathbf{N}\). Then p is not a zero divisor in A, and the set of a in A such that \(\theta(a) \equiv a^{p} \mod p\). A is a subring of A containing the X\_n and the Y\_n, hence equal to A.

By prop. 2, a) and c) of no. 2, there exist elements \(\mathbf{S} = (\mathrm{S}_{n})_{n\in\mathbb{N}}\) , \(\mathbf{P} = (\mathrm{P}_{n})_{n\in\mathbb{N}}\) , \(\mathbf{I} = (\mathrm{I}_{n})_{n\in\mathbb{N}}\) and \(\mathbf{F} = (\mathrm{F}_{n})_{n\in\mathbb{N}}\) of \(\mathrm{A}^{\mathbf{N}}\) characterized respectively by the equalities

\[
\left\{ \begin{array}{l} \Phi_ {\mathbf {A}} (\mathbf {S}) = \Phi_ {\mathbf {A}} (\mathbf {X}) + \Phi_ {\mathbf {A}} (\mathbf {Y}) \\ \Phi_ {\mathbf {A}} (\mathbf {P}) = \Phi_ {\mathbf {A}} (\mathbf {X}) \Phi_ {\mathbf {A}} (\mathbf {Y}) \\ \Phi_ {\mathbf {A}} (\mathbf {I}) = - \Phi_ {\mathbf {A}} (\mathbf {X}) \\ \Phi_ {\mathbf {A}} (\mathbf {F}) = f _ {\mathbf {A}} \big (\Phi_ {\mathbf {A}} (\mathbf {X}) \big). \end{array} \right.\tag{11}
\]

The elements \(S_{n}\), \(P_{n}\), \(I_{n}\), and \(F_{n}\) of A are therefore characterized by the following formulas (where \(n\) ranges over \(\mathbf{N}\)):

\begin{enumerate}
\def\labelenumi{(\arabic{enumi})}
\setcounter{enumi}{11}
\tightlist
\item
\end{enumerate}

\[
\Phi_ {n} \left(\mathrm{S} _ {0}, \dots , \mathrm{S} _ {n}\right) = \Phi_ {n} \left(\mathrm{X} _ {0}, \dots , \mathrm{X} _ {n}\right) + \Phi_ {n} \left(\mathrm{Y} _ {0}, \dots , \mathrm{Y} _ {n}\right),\tag{13}
\]

\[
\Phi_ {n} \left(\mathrm{P} _ {0}, \dots , \mathrm{P} _ {n}\right) = \Phi_ {n} \left(\mathrm{X} _ {0}, \dots , \mathrm{X} _ {n}\right) \Phi_ {n} \left(\mathrm{Y} _ {0}, \dots , \mathrm{Y} _ {n}\right),\tag{14}
\]

\[
\Phi_ {n} (\mathrm{I} _ {0}, \dots , \mathrm{I} _ {n}) = - \Phi_ {n} (\mathrm{X} _ {0}, \dots , \mathrm{X} _ {n}),\tag{15}
\]

\[
\Phi_ {n} \left(\mathrm{F} _ {0}, \dots , \mathrm{F} _ {n}\right) = \Phi_ {n + 1} \left(\mathrm{X} _ {0}, \dots , \mathrm{X} _ {n + 1}\right).
\]

Let us assign \(X_{n}\) and \(Y_{n}\) the weight \(p^{n}\) for every \(n \in \mathbf{N}\). From the remark of no. 2, we deduce the following assertions:

\begin{enumerate}
\def\labelenumi{\alph{enumi})}
\item
  We have \(\mathbf{S}_n\in \mathbf{Z}[\mathbf{X}_0,\dots ,\mathbf{X}_n,\mathbf{Y}_0,\dots ,\mathbf{Y}_n]\), and \(\mathbf{S}_n\) is isobaric of weight \(p^n\).
\item
  We have \(\mathbf{P}_n\in \mathbf{Z}[\mathrm{X}_0,\dots ,\mathrm{X}_n,\mathrm{Y}_0,\dots ,\mathrm{Y}_n]\), and \(\mathbf{P}_n\) is isobaric of weight \(p^n\) in each of the families \((\mathrm{X}_0,\dots ,\mathrm{X}_n)\) and \((\mathrm{Y}_0,\dots ,\mathrm{Y}_n)\).
\item
  We have \(\mathbf{I}_n\in \mathbf{Z}[\mathrm{X}_0,\dots ,\mathrm{X}_n]\), and \(\mathbf{I}_n\) is isobaric of weight \(p^n\).
\item
  We have \(\mathbf{F}_n\in \mathbf{Z}[\mathbf{X}_0,\dots ,\mathbf{X}_{n + 1}]\), and \(\mathbf{F}_n\) is isobaric of weight \(p^{n + 1}\).
\end{enumerate}

Formula (2) allows in practice the determination of the polynomials \(S_{n}\) , \(P_{n}\) , \(I_{n}\) and \(F_{n}\) step by step.

Examples. -- 1) We have

\[
\mathbf {S _ {0}} = \mathbf {X _ {0}} + \mathbf {Y _ {0}}
\]

\[
\mathrm{S} _ {1} = \mathrm{X} _ {1} + \mathrm{Y} _ {1} - \sum_ {i = 1} ^ {p - 1} \frac {1}{p} \binom {p} {i} \mathrm{X} _ {0} ^ {i} \mathrm{Y} _ {0} ^ {p - i}.
\]

Additionally, \(S_{n} - X_{n} - Y_{n}\) belongs to the ring \(Z[X_{0}, ..., X_{n-1}, Y_{0}, ..., Y_{n-1}]\) . 2) We have

\[
\begin{array}{l} \mathbf {P} _ {0} = \mathbf {X} _ {0} \mathbf {Y} _ {0} \\ \mathbf {P} _ {1} = p \mathbf {X} _ {1} \mathbf {Y} _ {1} + \mathbf {X} _ {0} ^ {p} \mathbf {Y} _ {1} + \mathbf {X} _ {1} \mathbf {Y} _ {0} ^ {p}. \end{array}
\]

\begin{enumerate}
\def\labelenumi{\arabic{enumi})}
\setcounter{enumi}{2}
\tightlist
\item
  When \(p \neq 2\) , one has \(I_n = -X_n\) . For \(p = 2\) , one has
\end{enumerate}

\[
\mathrm{I} _ {0} = - \mathrm{X} _ {0}
\]

\[
\mathrm{I} _ {1} = - \left(\mathrm{X} _ {0} ^ {2} + \mathrm{X} _ {1}\right)
\]

\[
\mathrm{I} _ {2} = - \mathrm{X} _ {0} ^ {4} - \mathrm{X} _ {0} ^ {2} \mathrm{X} _ {1} - \mathrm{X} _ {1} ^ {2} - \mathrm{X} _ {2}.
\]

\begin{enumerate}
\def\labelenumi{\arabic{enumi})}
\setcounter{enumi}{3}
\tightlist
\item
  We have
\end{enumerate}

\[
\mathbf {F} _ {0} = \mathbf {X} _ {0} ^ {p} + p \mathbf {X} _ {1}
\]

\[
\mathrm{F} _ {1} = \mathrm{X} _ {1} ^ {p} + p \mathrm{X} _ {2} - \sum_ {i = 0} ^ {p - 1} \binom {p} {i} p ^ {p - i - 1} \mathrm{X} _ {0} ^ {p i} \mathrm{X} _ {1} ^ {p - i}.
\]

Since \(\Phi_n(F_0, \ldots, F_n) \equiv \Phi_n(X_0^p, \ldots, X_n^p) \pmod{p^{n+1}A}\) for all \(n \in \mathbf{N}\) (formulas (2) and (15)), it follows from Proposition 1(b) that \(F_n \equiv X_n^p \pmod{pA}\) for all \(n \in \mathbf{N}\).

Remark. --- Let J be the set of integers \(j \geqslant 1\) . For every element \(j\) of J, define the polynomial \(\varphi_j\) of \(\mathbf{Z}[(\mathbf{X}_j)_{j \in \mathbf{J}}]\) by the formula

\[
\varphi_ {j} = \sum_ {d} d X _ {d} ^ {j / d},
\]

where the sum is over the elements of J that divide j. For every integer \(n \geqslant 0\) , one has

\[
\varphi_ {p ^ {n}} = \Phi_ {n} (\mathrm{X} _ {p ^ {0}},..., \mathrm{X} _ {p ^ {n}}).
\]

For every ring A and every element \(m\) of \(\mathbf{J}\), we denote by \(\varphi_{m}\) the map from \(A^{\mathbf{J}}\) into A that sends \((a_{j})_{j\in\mathbf{J}}\) to \(\varphi_{m}((a_{j})_{j\in\mathbf{J}})\); we denote by \(\varphi_{A}\), or simply \(\varphi\), the map from \(A^{\mathbf{J}}\) into itself that sends \(a = (a_{j})_{j\in\mathbf{J}}\) to \((\varphi_{m}(a))_{m\in\mathbf{J}}\).

Let \(\mathcal{A} = \mathbf{Z}[(X_j)_{j\in\mathbf{J}},(Y_j)_{j\in\mathbf{J}}]\) be the ring of polynomials with integer coefficients in the two families of indeterminates \(\mathbf{X} = (X_j)_{j\in\mathbf{J}}\) and \(\mathbf{Y} = (Y_j)_{j\in\mathbf{J}}\). One can show (p. 51, exerc. 34) that there exist in \(\mathcal{A}\) elements

\[
\boldsymbol {s} = (s _ {j}) _ {j \in \mathbf {J}}, \boldsymbol {p} = (p _ {j}) _ {j \in \mathbf {J}} \quad \text{and} \quad \boldsymbol {i} = (i _ {j}) _ {j \in \mathbf {J}},
\]

characterized by the following equalities:

\[
\begin{array}{l} \varphi_ {\mathcal {A}} (s) = \varphi_ {\mathcal {A}} (\mathbf {X}) + \varphi_ {\mathcal {A}} (\mathbf {Y}) \\ \varphi_ {\mathcal {A}} (p) = \varphi_ {\mathcal {A}} (\mathbf {X}) \varphi_ {\mathcal {A}} (\mathbf {Y}) \\ \varphi_ {\mathcal {A}} (i) = - \varphi_ {\mathcal {A}} (\mathbf {X}). \end{array}
\]

